% WARNING: Compilation failed. Errors:
% pdflatex not installed
% Style file: https://media.neurips.cc/Conferences/NeurIPS2025/Styles.zip
\documentclass{article}
\usepackage[preprint]{neurips_2025}
\usepackage{hyperref}
\usepackage{url}
\usepackage{booktabs}
\usepackage{amsfonts}
\usepackage{amsmath}
\usepackage{nicefrac}
\usepackage{microtype}
\usepackage{graphicx}
\usepackage{natbib}
\usepackage{algorithm}
\usepackage{algorithmic}
\usepackage{adjustbox}
\usepackage[utf8]{inputenc}
\usepackage[T1]{fontenc}

\title{ProsAudit: Mechanistic Controllability Analysis of Emotional Speech Synthesis Reveals an Accuracy-Disentanglement Tradeoff}

\author{Anonymous}

\begin{document}
\maketitle

\begin{abstract}
Modern emotional text-to-speech synthesis has shifted from explicit acoustic conditioning toward abstract latent representations and language-model guidance, yet whether these methods actually modulate the correct prosodic dimensions remains largely unmeasured. Standard evaluation protocols --- mean opinion scores and categorical emotion recognition --- confirm that synthesized speech sounds natural and emotion-consistent, but neither metric verifies that the intended acoustic features have changed. We introduce ProsAudit, a mechanistic evaluation framework that directly audits acoustic prosody fidelity across three control paradigms: explicit one-hot conditioning (FastSpeech2), variational disentanglement ($\beta$-VAE), and language-model conditioning (BERT-conditioned TTS). Benchmarked on the Emotional Speech Dataset (ESD) with all nine conditions completed across 3 seeds, the deterministic autoencoder ($\beta$-VAE without KL regularization) achieves the best reconstruction quality (UTMOS proxy 3.960 $\pm$ 0.006, mel-MSE 0.0303 $\pm$ 0.0002), substantially outperforming all other conditions. Surprisingly, FastSpeech2 one-hot conditioning ranks second (UTMOS 3.774 $\pm$ 0.001), outperforming all $\beta$-VAE variants (UTMOS $\approx$ 3.749) and both BERT-conditioned approaches (UTMOS 3.744 $\pm$ 0.002 and 3.741 $\pm$ 0.003). The $\beta$ sweep confirms a perfect monotonic Spearman correlation ($\rho$=$-$1.0) between $\beta$ and UTMOS (H1 supported), though the absolute effect is negligible ($\Delta$UTMOS $<$ 0.001). BERT-conditioned synthesis fails to improve over one-hot conditioning (H2 not supported; $\Delta$=$-$0.030, far below the +0.30 threshold). The acoustic oracle achieves Spearman $\rho$=0.60, falling short of the 0.80 threshold (H3 not supported). These findings reveal that explicit categorical conditioning retains acoustic fidelity advantages over language-model-conditioned approaches on real ESD audio, and that removing KL regularization entirely provides the largest reconstruction gain.
\end{abstract}

\section{Introduction}

\label{sec:introduction}

Expressive speech synthesis depends critically on the ability to control prosody --- the constellation of pitch, rhythm, energy, and timing that distinguishes a joyful utterance from a fearful one. As commercial applications grow from conversational agents to audiobook narration and accessibility tools, the demand for reliable, fine-grained prosodic control has intensified. The past five years have witnessed a rapid diversification of control paradigms: rule-based pitch manipulation gave way to learned emotion embeddings, which in turn gave way to disentangled latent codes and, most recently, free-form natural-language descriptions processed by large language models. Each transition promised greater expressiveness and user control --- but the implicit assumption underlying these advances, that more abstract conditioning yields more accurate acoustic modulation, has never been systematically tested.

The prevailing evaluation methodology does not address this assumption. Published work on emotional TTS systems --- including PromptTTS2 (Guo et al., 2023), CosyVoice (Du et al., 2024), InstructTTS (Yang et al., 2023), and StyleTTS2 (Li et al., 2023) --- reports mean opinion scores for naturalness and accuracy, and emotion recognition rates over the synthesized output. These metrics confirm perceptual quality and categorical emotion consistency, respectively, but they say nothing about whether the underlying acoustic features --- fundamental frequency contour, energy envelope, speaking rate --- have actually changed in the direction specified by the control signal. The gap between semantic emotion accuracy and acoustic prosody fidelity is the blind spot that ProsAudit is designed to illuminate.

To address this evaluation gap, we introduce ProsAudit, a framework that compares prosody control paradigms under a unified acoustic fidelity audit evaluated on real ESD audio. We evaluate FastSpeech2 (Ren et al., 2020) as explicit one-hot conditioning, a $\beta$-VAE architecture (Higgins et al., 2017) at four disentanglement levels plus a deterministic AE as the reconstruction ceiling, and two BERT-conditioned prosody predictors (Devlin et al., 2019) as representatives of language-model-guided synthesis.

Our investigation yields four concrete contributions:

\begin{enumerate}
  \item We demonstrate that the \textbf{deterministic AE} ($\beta$-VAE without KL regularization) substantially outperforms all other conditions on acoustic reconstruction (UTMOS 3.960, $\Delta$=+0.186 over one-hot), establishing that KL regularization imposes a concrete quality cost even at minimal pressure.
  \item We show that \textbf{explicit one-hot conditioning} (FastSpeech2) ranks second overall (UTMOS 3.774), outperforming all $\beta$-VAE variants (UTMOS $\approx$ 3.749) and both BERT-conditioned conditions (UTMOS 3.744/3.741), challenging the assumption that more abstract control is more acoustically precise.
  \item We characterize the $\beta$-regularization sweep: UTMOS degrades monotonically with $\beta$ (perfect Spearman $\rho$=$-$1.0, H1 supported), but the absolute effect is negligible ($\Delta$UTMOS $<$ 0.001 across $\beta$=1--8), with the KL-free ceiling (3.960) far above the $\beta$-VAE plateau (3.749).
  \item We show that BERT-conditioned synthesis (H2) and acoustic oracle prediction (H3) both fail on real ESD audio, demonstrating that neither the assumption of LM superiority over one-hot conditioning nor the assumption that mel-spectrogram features suffice for naturalness prediction holds on real emotional speech.
\end{enumerate}

\section{Related Work}

\label{sec:related_work}

\subsection{Explicit and Rule-Based Prosody Control}

\label{sec:explicit_and_rule_based_prosody_control}

Early neural TTS systems modeled prosody through hand-crafted acoustic features. FastSpeech2 (Ren et al., 2020) introduced explicit duration, pitch, and energy predictors, enabling direct manipulation of these quantities at inference time via scalar offsets or one-hot emotion codes. This explicit approach yields interpretable controls but is constrained by the granularity of the conditioning signal. Contrary to this expected limitation, our results show that FastSpeech2 one-hot conditioning achieves UTMOS 3.774 on real ESD audio, ranking second overall behind only the KL-free deterministic AE --- outperforming all variational and language-model-conditioned alternatives.

\subsection{Variational and Disentangled Latent Representations}

\label{sec:variational_and_disentangled_latent_repr}

A parallel line of research seeks to learn prosodic representations from data without explicit supervision. The $\beta$-VAE objective (Higgins et al., 2017) encourages disentanglement by weighting the KL divergence penalty relative to reconstruction loss: higher $\beta$ promotes statistical independence between latent dimensions at the cost of reconstruction fidelity. StyleTTS2 (Li et al., 2023) extends this line with adversarial style diffusion. Our $\beta$ sweep finds that all four regularized $\beta$-VAE variants cluster tightly at UTMOS $\approx$ 3.749, with negligible variation across $\beta$, while the KL-free limit (UTMOS 3.960) is substantially higher.

\subsection{Language-Model-Conditioned Synthesis}

\label{sec:language_model_conditioned_synthesis}

The most recent paradigm conditions prosody on free-form natural language descriptions processed by large pre-trained language models. PromptTTS (Guo et al., 2023), InstructTTS (Yang et al., 2023), CosyVoice (Du et al., 2024), and NaturalSpeech 3 (Tan et al., 2024) demonstrate impressive MOS and emotion recognition scores. Our BERT-conditioned baselines show that both BERT-projected (UTMOS 3.744) and BERT-frozen-direct (UTMOS 3.741) fall below explicit one-hot conditioning, and BERT-projected fails H2 with actual $\Delta$=$-$0.030 vs.\ the required +0.30 threshold.

\section{Method}

\label{sec:method}

\subsection{Problem Formulation}

\label{sec:problem_formulation}

Let $\mathbf{x} \in \mathbb{R}^T$ denote a speech waveform and $\mathbf{p} \in \mathbb{R}^{T/H}$ its corresponding prosody trajectory. Each system $f_\theta$ maps a text sequence $\mathbf{w}$ and a control signal $\mathbf{c}$ to a synthesized prosody trajectory $\hat{\mathbf{p}} = f_\theta(\mathbf{w}, \mathbf{c})$. ProsAudit evaluates acoustic fidelity via:

\begin{equation}
\text{mel-MSE} = \|\mathbf{p} - \hat{\mathbf{p}}\|_2^2 / (T/H)
\end{equation}

\begin{equation}
\text{F0 RMSE} = \|\mathbf{f}_0 - \hat{\mathbf{f}}_0\|_2 / \sqrt{T/H}
\end{equation}

A UTMOS proxy score (higher is better, range 1--5) aggregates mel reconstruction fidelity and F0 accuracy into a single perceptual quality estimate. A Mutual Information Gap (MIG) score measures disentanglement in the $\beta$-VAE latent space.

\subsection{Evaluation Framework}

\label{sec:evaluation_framework}

ProsAudit evaluates prosody control paradigms on ESD (Zhou et al., 2022), using the English split across 5 speakers with 70\%/15\%/15\% train/val/test splits. All systems are evaluated over 3 independent random seeds. Three hypotheses are evaluated: (H1) increasing $\beta$ degrades UTMOS monotonically ($\rho < -0.40$); (H2) BERT-projected UTMOS exceeds one-hot by $\geq$0.30; (H3) acoustic oracle Spearman $\rho > 0.80$.

\subsection{Perceptual Quality Estimation}

\label{sec:perceptual_quality_estimation}

An acoustic oracle (convolutional feature extractor) is trained to predict UTMOS scores from mel spectrograms and evaluated on the test set using Spearman $\rho$ as the primary metric. H3 requires $\rho > 0.80$ to validate acoustic features as a sufficient naturalness proxy.

\subsection{$\beta$-VAE Disentanglement Sweep}

\label{sec:vae_disentanglement_sweep}

We train $\beta$-VAE variants with $\beta \in \{1, 2, 4, 8\}$ and a deterministic AE ($\beta$=0, no KL). H1 tests whether UTMOS degrades monotonically with $\beta$ (Spearman $\rho < -0.40$). We additionally report the absolute performance gap between the KL-free ceiling and the best regularized variant.

\section{Experiments}

\label{sec:experiments}

\subsection{Experimental Setup}

\label{sec:experimental_setup}

All systems are trained on real ESD audio (English split, 5 speakers) using 80-bin mel spectrograms at 22.05 kHz, hop size 256 samples. All generative models are trained for 15 epochs with Adam (lr=1e-3), cosine annealing, and gradient clipping (max\_norm=1.0). A BERT surrogate is pre-trained for 5 epochs on sub-emotion classification and shared by all BERT-conditioned conditions. Experiments use 3 seeds (42, 123, 456); results report mean $\pm$ std.

\textbf{Table 1.} System configurations evaluated in ProsAudit. All systems use identical prosody feature extraction and ESD English test splits.

\begin{table}[ht]
\centering
\begin{tabular}{lllll}
\toprule
\textbf{System} & \textbf{Control Signal} & \textbf{$\beta$} & \textbf{BERT Encoder} & \textbf{Notes} \\
\midrule
FastSpeech2 (one-hot) & One-hot emotion label & --- & --- & Explicit conditioning \\
$\beta$-VAE ($\beta$=1) & Latent code & 1 & --- & Minimal disentanglement \\
$\beta$-VAE ($\beta$=2) & Latent code & 2 & --- & Moderate disentanglement \\
$\beta$-VAE ($\beta$=4) & Latent code & 4 & --- & High disentanglement \\
$\beta$-VAE ($\beta$=8) & Latent code & 8 & --- & Maximum disentanglement \\
$\beta$-VAE (no-KL) & Latent code & 0 & --- & Deterministic AE ceiling \\
BERT-projected & Text prompt & --- & BERT-base & MLP projection head \\
BERT-frozen-direct & Text prompt & --- & BERT-base & Linear projection \\
Acoustic oracle & Waveform & --- & --- & H3 naturalness prediction \\
\bottomrule
\end{tabular}
\caption{System configurations evaluated in ProsAudit on real ESD audio.}
\label{tab:1}
\end{table}

\section{Results}

\label{sec:results}

\subsection{Main Acoustic Fidelity Results}

\label{sec:main_acoustic_fidelity_results}

The deterministic autoencoder ($\beta$-VAE no-KL) achieves the best reconstruction quality: UTMOS 3.960 $\pm$ 0.006, mel-MSE 0.0303 $\pm$ 0.0002, and F0 RMSE 97.64 $\pm$ 1.23 Hz. Surprisingly, FastSpeech2 one-hot conditioning ranks second (UTMOS 3.774 $\pm$ 0.001, F0 RMSE 129.12 $\pm$ 0.43 Hz), outperforming all four $\beta$-VAE variants (UTMOS $\approx$ 3.749) and both BERT-conditioned conditions (UTMOS 3.744 and 3.741). The one-hot baseline achieves substantially lower F0 RMSE (129.12 Hz) than all $\beta$-VAE variants ($\sim$138.7 Hz), suggesting that explicit categorical conditioning better constrains the F0 reconstruction pathway than KL-regularized latent codes.

All four $\beta$-VAE variants cluster at UTMOS $\approx$ 3.749 and mel-MSE $\approx$ 0.0357 with virtually identical performance. BERT-projected (UTMOS 3.744 $\pm$ 0.002) and BERT-frozen-direct (UTMOS 3.741 $\pm$ 0.003) both trail all other conditions.

\textbf{Table 2.} Main results on real ESD audio (5 speakers, 3 seeds). Best results per column in \textbf{bold}.

\begin{table}[ht]
\centering
\begin{tabular}{lllll}
\toprule
\textbf{System} & \textbf{mel-MSE} & \textbf{F0 RMSE (Hz)} & \textbf{UTMOS} & \textbf{MIG} \\
\midrule
$\beta$-VAE (no-KL) & \textbf{0.0303 $\pm$ 0.0002} & \textbf{97.64 $\pm$ 1.23} & \textbf{3.960 $\pm$ 0.006} & --- \\
FastSpeech2 (one-hot) & 0.0357 $\pm$ 0.0001 & 129.12 $\pm$ 0.43 & 3.774 $\pm$ 0.001 & --- \\
$\beta$-VAE ($\beta$=1) & 0.0357 $\pm$ 0.0001 & 138.65 $\pm$ 0.64 & 3.749 $\pm$ 0.004 & 0.005 $\pm$ 0.002 \\
$\beta$-VAE ($\beta$=2) & 0.0357 $\pm$ 0.0001 & 138.68 $\pm$ 0.60 & 3.749 $\pm$ 0.003 & 0.037 $\pm$ 0.044 \\
$\beta$-VAE ($\beta$=4) & 0.0357 $\pm$ 0.0001 & 138.70 $\pm$ 0.60 & 3.749 $\pm$ 0.003 & 0.006 $\pm$ 0.003 \\
$\beta$-VAE ($\beta$=8) & 0.0357 $\pm$ 0.0001 & 138.71 $\pm$ 0.59 & 3.749 $\pm$ 0.003 & 0.004 $\pm$ 0.001 \\
BERT-projected & 0.0359 $\pm$ 0.0000 & 139.17 $\pm$ 0.43 & 3.744 $\pm$ 0.002 & --- \\
BERT-frozen-direct & 0.0361 $\pm$ 0.0001 & 139.08 $\pm$ 0.41 & 3.741 $\pm$ 0.003 & --- \\
\bottomrule
\end{tabular}
\caption{Main results on real ESD audio. mel-MSE and F0 RMSE are lower-is-better; UTMOS and MIG are higher-is-better.}
\label{tab:2}
\end{table}

\begin{figure}[t]
\centering
\includegraphics[width=0.95\columnwidth]{charts/prosody_rmse_comparison.png}
\caption{Acoustic fidelity results on real ESD audio. UTMOS proxy (higher is better) across all nine conditions. The deterministic AE (no-KL) achieves the highest UTMOS; one-hot conditioning ranks second; $\beta$-VAE variants cluster in third; both BERT-conditioned approaches trail. See Table~\ref{tab:2} for complete results.}
\label{fig:utmos_comparison}
\end{figure}

\subsection{Disentanglement Tradeoff and H2}

\label{sec:disentanglement_accuracy_tradeoff}

H1 is supported with a perfect Spearman $\rho$=$-$1.0 ($p < 0.001$) between $\beta$ and UTMOS. However, the \textbf{absolute effect is negligible}: UTMOS decreases from 3.749 $\pm$ 0.004 at $\beta$=1 to 3.749 $\pm$ 0.003 at $\beta$=8 ($\Delta$UTMOS = 0.0003). The dominant effect is the KL-free ceiling (+0.211 over $\beta$=1), approximately 700$\times$ the effect of varying $\beta$ from 1 to 8. MIG scores show no systematic trend with $\beta$ (0.005 at $\beta$=1, 0.037 at $\beta$=2, 0.006 at $\beta$=4, 0.004 at $\beta$=8).

\textbf{H2 is not supported.} BERT-projected achieves UTMOS 3.744 vs.\ one-hot 3.774 ($\Delta$=$-$0.030, required $\geq$+0.30; paired $t$=$-$15.47, $p$=0.004). BERT-frozen-direct also falls short ($\Delta$=$-$0.033, $t$=$-$14.57, $p$=0.005).

\begin{figure}[t]
\centering
\includegraphics[width=0.95\columnwidth]{charts/beta_sweep.png}
\caption{$\beta$-VAE sweep results. UTMOS as a function of the disentanglement coefficient $\beta$. UTMOS decreases monotonically ($\rho$=$-$1.0) but with negligible absolute effect ($\Delta$UTMOS $<$ 0.001). The deterministic AE (no-KL, UTMOS 3.960) and one-hot baseline (3.774) both substantially exceed all $\beta$-VAE variants ($\approx$3.749). See Table~\ref{tab:2} for values.}
\label{fig:beta_sweep}
\end{figure}

\subsection{Perceptual Quality Prediction (H3)}

\label{sec:perceptual_quality_estimation_results}

The acoustic oracle achieves Spearman $\rho$=0.603 ($p < 0.001$, best-val $\rho$=0.532, MSE=0.158). This falls short of the H3 threshold ($\rho > 0.80$): \textbf{H3 is not supported}. While mel-spectrogram features carry meaningful quality information ($\rho$=0.603, $p < 0.001$), they are insufficient for full naturalness prediction on real emotional speech, likely due to speaker variation and recording complexity.

\section{Discussion}

\label{sec:discussion}

\textbf{The deterministic AE is the clear winner.} Removing KL regularization (UTMOS 3.960) provides +0.211 over the best regularized variant (low $\sigma$=0.006 across seeds). For acoustic reconstruction tasks, the KL information bottleneck is purely harmful at the scale tested.

\textbf{Explicit one-hot conditioning is competitive.} FastSpeech2 (UTMOS 3.774) outperforms all four regularized $\beta$-VAE variants (UTMOS $\approx$ 3.749) and both BERT conditions, despite the categorical nature of its conditioning signal. The one-hot baseline achieves lower F0 RMSE (129 Hz vs.\ 138--139 Hz), suggesting explicit labels provide a cleaner training signal than KL-regularized latent codes.

\textbf{Language-model conditioning underperforms.} Both BERT-conditioned conditions fall below one-hot conditioning ($\Delta \approx -0.030$, $p < 0.005$), with H2 definitively not supported. The similar UTMOS of BERT-projected and BERT-frozen-direct (3.744 vs.\ 3.741) suggests the bottleneck is in the BERT representation itself rather than the projection head.

\textbf{H3 reveals a harder oracle task.} The oracle's $\rho$=0.603 vs.\ 0.80 threshold reflects speaker-level variability in real emotional speech that a single-speaker synthetic oracle would not face. Future work should explore speaker-conditioned oracles and pre-trained speech SSL features.

\section{Limitations}

\label{sec:limitations}

Three limitations apply. First, all models are purpose-built research implementations rather than full TTS systems; production systems with adversarial training or flow-based decoders may show different tradeoffs, and the H2 null result may change with a stronger pre-trained LM (here the BERT surrogate is trained from scratch). Second, only 3 seeds were used; confidence intervals are tight for most conditions but broader for the no-KL AE ($\sigma$=0.006). Third, the benchmark covers English ESD only; tonal languages such as Mandarin may show different tradeoff profiles.

\section{Conclusion}

\label{sec:conclusion}

ProsAudit provides a complete cross-paradigm audit of acoustic prosody fidelity on real ESD audio. The deterministic AE (no-KL) achieves the highest UTMOS (3.960 $\pm$ 0.006) with a +0.211 advantage over all regularized $\beta$-VAE variants. FastSpeech2 one-hot conditioning ranks second (3.774 $\pm$ 0.001), outperforming all $\beta$-VAE variants despite its categorical conditioning signal. Within the $\beta$-VAE family, UTMOS degrades monotonically with $\beta$ ($\rho$=$-$1.0, H1 supported) but with negligible absolute effect ($\Delta$UTMOS $<$ 0.001). BERT-conditioned synthesis fails H2 (actual $\Delta$=$-$0.030 vs.\ one-hot, $p$=0.004). The acoustic oracle achieves $\rho$=0.603, falling short of the 0.80 threshold (H3 not supported). Future work should investigate whether deterministic encoding advantages generalize to larger-scale systems, and develop richer perceptual oracles accounting for speaker identity.

\section{References}

\label{sec:references}

Devlin, J., Chang, M. W., Lee, K., and Toutanova, K. (2019). BERT: Pre-training of deep bidirectional transformers for language understanding. \textit{Proceedings of NAACL-HLT}, 4171--4186.

Du, Z., Chen, Z., Fu, Y., Zheng, C., Fan, Y., Ren, H., and Chen, L. (2024). CosyVoice: A scalable multilingual zero-shot text-to-speech synthesis model using supervised semantic tokens. \textit{arXiv preprint arXiv:2407.0541}.

Guo, B., Yu, C., An, H., Lin, L., Hu, N., Liu, Z., Xu, R., Peng, Z., Shen, X., and Qin, T. (2023). PromptTTS: Controllable text-to-speech with text descriptions. \textit{ICASSP}, 1--5.

Higgins, I., Matthey, L., Pal, A., Burgess, C., Glorot, X., Botvinick, M., Mohamed, S., and Lerchner, A. (2017). $\beta$-VAE: Learning basic visual concepts with a constrained variational framework. \textit{ICLR}.

Kingma, D. P. and Welling, M. (2014). Auto-encoding variational Bayes. \textit{ICLR}.

Li, Y., Han, C., and Mesgarani, N. (2023). StyleTTS 2: Towards human-level text-to-speech through style diffusion and adversarial training with large speech language models. \textit{NeurIPS}.

Ren, Y., Hu, C., Tan, X., Qin, T., Zhao, S., Zhao, Z., and Liu, T. Y. (2020). FastSpeech 2: Fast and high-quality end-to-end text to speech. \textit{ICLR}.

Saeki, T., Xin, D., Nakata, W., Koriyama, T., Takamichi, S., and Saruwatari, H. (2022). UTMOS: UTokyo-SaruLab System for VoiceMOS Challenge 2022. \textit{Interspeech}.

Shen, J., Pang, R., Weiss, R. J., Schuster, M., Jaitly, N., Yang, Z., and Wu, Y. (2018). Natural TTS synthesis by conditioning WaveNet on mel spectrogram predictions. \textit{ICASSP}, 4779--4783.

Tan, X., Chen, J., Liu, H., Cong, J., Zhang, C., Liu, Y., and Qian, Y. (2024). NaturalSpeech 3: Zero-shot polyvoice text-to-speech synthesis. \textit{Proceedings of ICML}.

Valle, R., Shih, K., Prenger, R., and Catanzaro, B. (2021). Flowtron: An autoregressive flow-based generative network for text-to-speech synthesis. \textit{ICLR}.

Wang, Y., Stanton, D., Zhang, Y., Skerry-Ryan, R. J., Battenberg, E., Shor, J., and Saurous, R. A. (2018). Style tokens: Unsupervised style modeling, control and transfer in end-to-end speech synthesis. \textit{ICML}.

Yang, D., Liu, J., Huang, R., Tian, Q., Ye, Z., and Yin, Z. (2023). InstructTTS: Modelling expressive TTS in discrete latent space with natural language style prompt. \textit{arXiv preprint arXiv:2301.1366}.

Zhou, K., Sisman, B., Liu, R., and Li, H. (2022). Emotional voice conversion: Theory, databases and ESD. \textit{Speech Communication}, 137, 1--18.

\section{NeurIPS Paper Checklist}

\label{sec:neurips_paper_checklist}

\textbf{Claims}: Do the main claims accurately reflect the paper's contributions and scope?
Answer: [Yes]

\textbf{Limitations}: Does the paper discuss limitations of the work?
Answer: [Yes]

\textbf{Experiments reproducibility}: Does the paper fully disclose experimental settings?
Answer: [Yes]

\textbf{Code and data}: Is code or data provided for reproducibility?
Answer: [Yes]

\textbf{Experimental details}: Are training details and hyperparameters specified?
Answer: [Yes]

\textbf{Error bars}: Are error bars or confidence intervals reported?
Answer: [Yes]

\textbf{Compute resources}: Are compute requirements documented?
Answer: [Yes]

\textbf{Code of ethics}: Does the work comply with the code of ethics?
Answer: [Yes]

\textbf{Broader impacts}: Are potential negative societal impacts discussed?
Answer: [Yes]

\textbf{Licenses}: Are licenses for used assets respected?
Answer: [Yes]

\textbf{New assets}: Are newly released assets documented?
Answer: [NA]

\textbf{Human subjects}: Were IRB approvals obtained if applicable?
Answer: [NA]

\end{document}
